\section{Elementy realizacji zadania}
\label{sec:realizacja}

\subsection{Implementacja potoku}

Część praktyczną zrealizowano jako modułowy potok w języku Python (katalog
\texttt{src/}). Najważniejsze komponenty to:
\begin{itemize}
  \item \texttt{data/} --- generator danych syntetycznych, wczytywanie realnego
        zbioru CIC-IDS2018 oraz przetwarzanie wstępne (czyszczenie, kodowanie,
        standaryzacja, stratyfikowany podział);
  \item \texttt{models/classical.py} --- las losowy i regresja logistyczna
        (\texttt{scikit-learn});
  \item \texttt{models/deep.py} --- architektury MLP, CNN 1D i LSTM (PyTorch) z
        wspólną pętlą treningową, wczesnym zatrzymaniem i ważeniem klas;
  \item \texttt{evaluate.py}, \texttt{plots.py}, \texttt{tables.py} --- obliczanie
        metryk, generowanie rysunków oraz tabel i makr \LaTeX;
  \item \texttt{run\_all.py} --- punkt wejścia uruchamiający cały eksperyment.
\end{itemize}
Wszystkie wyniki liczbowe, tabele i rysunki w tym rozdziale zostały
wygenerowane automatycznie przez ten kod, dzięki czemu treść sprawozdania jest
ściśle spójna z implementacją.

\subsection{Skuteczność modeli}

Tabela~\ref{tab:metrics} zestawia metryki skuteczności wszystkich modeli na
zbiorze testowym (wartości najlepsze w każdej kolumnie wyróżniono pogrubieniem).

\begin{table}[H]
\centering
\caption{Porównanie skuteczności modeli na zbiorze testowym.}
\label{tab:metrics}
\small
\begin{tabular}{lrrrrrr}
\toprule
Model & Accuracy & Precision & Recall & F1 & ROC-AUC & PR-AUC \\
\midrule
Logistic Regression & 0.886 & 0.699 & 0.890 & 0.783 & 0.945 & 0.875 \\
Random Forest & 0.961 & 0.932 & \textbf{0.893} & 0.912 & 0.968 & \textbf{0.953} \\
MLP & 0.954 & 0.912 & 0.884 & 0.898 & 0.967 & 0.946 \\
CNN 1D & \textbf{0.961} & \textbf{0.942} & 0.885 & \textbf{0.913} & \textbf{0.969} & 0.949 \\
LSTM & 0.947 & 0.895 & 0.870 & 0.882 & 0.943 & 0.913 \\
\bottomrule
\end{tabular}
\end{table}

Z zestawienia wynika kilka obserwacji. Po pierwsze, wszystkie metody bardziej
zaawansowane wyraźnie przewyższają prostą regresję logistyczną --- potwierdza to,
że atak zdradza dopiero \emph{kombinacja} cech, a nie pojedyncze wartości. Po
drugie, na samym szczycie panuje praktycznie remis: najlepszy model głęboki
\textbf{\resBestDeepModel} ($F_1 = \resBestDeepFOne$) i \textbf{las losowy}
($F_1 = \resRFFOne$, ROC-AUC $= \resRFAUC$) osiągają niemal identyczne wyniki ---
różnica jest na granicy szumu. Po trzecie, oba modele mają nieco inny
,,charakter'': las losowy uzyskuje najwyższą \emph{czułość} (najwięcej wykrytych
ataków) oraz najwyższy PR-AUC, podczas gdy CNN ma najwyższą \emph{precyzję}
(najmniej fałszywych alarmów). Warto też zauważyć, że najbardziej rozbudowany
model --- rekurencyjna sieć LSTM --- wypada najsłabiej spośród metod
nieliniowych, mimo największej złożoności.

Graficzne porównanie metryk przedstawia rys.~\ref{fig:metrics}, a krzywe ROC
i precyzja--czułość --- rys.~\ref{fig:roc} i~\ref{fig:pr}. Krzywe ROC dla
wszystkich modeli nieliniowych przebiegają blisko lewego górnego rogu, co
świadczy o dobrej separowalności klas.

\begin{figure}[H]
\centering
\includegraphics[width=0.82\textwidth]{metrics_comparison.png}
\caption{Porównanie metryk skuteczności modeli na zbiorze testowym.}
\label{fig:metrics}
\end{figure}

\begin{figure}[H]
\centering
\begin{minipage}{0.49\textwidth}
  \centering
  \includegraphics[width=\textwidth]{roc_curves.png}
  \caption{Krzywe ROC.}
  \label{fig:roc}
\end{minipage}\hfill
\begin{minipage}{0.49\textwidth}
  \centering
  \includegraphics[width=\textwidth]{pr_curves.png}
  \caption{Krzywe precyzja--czułość.}
  \label{fig:pr}
\end{minipage}
\end{figure}

\subsection{Analiza błędów}

Macierze pomyłek (rys.~\ref{fig:cm}) pozwalają zajrzeć ,,pod maskę'' i zobaczyć,
jakie błędy popełniają modele. Las losowy wychwytuje najwięcej ataków (najwyższa
czułość), czyli rzadziej przepuszcza zagrożenie --- co w kontekście ochrony
danych jest cenne, bo przeoczone naruszenie kosztuje najwięcej. Sieć CNN z kolei
nieco rzadziej podnosi fałszywy alarm (najwyższa precyzja), czyli mniej męczy
zespół bezpieczeństwa niepotrzebnymi powiadomieniami. W praktyce ,,punkt
równowagi'' między tymi błędami można regulować progiem decyzyjnym, dopasowując
go do polityki bezpieczeństwa danej organizacji.

\begin{figure}[H]
\centering
\includegraphics[width=0.82\textwidth]{confusion_matrices.png}
\caption{Macierze pomyłek na zbiorze testowym dla wszystkich modeli.}
\label{fig:cm}
\end{figure}

\subsection{Wydajność i koszt obliczeniowy}

Skuteczność to tylko jeden wymiar oceny; w środowisku chmurowym, gdzie modele
są często dotrenowywane i muszą klasyfikować ruch na bieżąco, kluczowy jest
także koszt. Tabela~\ref{tab:efficiency} zestawia czasy treningu i inferencji
oraz złożoność modeli.

\begin{table}[H]
\centering
\caption{Porównanie wydajności i złożoności modeli.}
\label{tab:efficiency}
\small
\begin{tabular}{lrrl}
\toprule
Model & Czas treningu [s] & Czas inferencji [s] & Złożoność modelu \\
\midrule
Logistic Regression & 0.78 & 0.0006 & 69 wsp. \\
Random Forest & 4.16 & 0.0801 & 300 drzew / 1\,550\,456 węzłów \\
MLP & 11.82 & 0.0281 & 17\,537 param. \\
CNN 1D & 171.74 & 0.7103 & 39\,553 param. \\
LSTM & 487.45 & 3.9718 & 134\,401 param. \\
\bottomrule
\end{tabular}
\end{table}

I tu pojawia się sedno sprawy. Las losowy uczy się w zaledwie
\resRFTrainTime~s, podczas gdy najlepszy model głęboki (\resBestDeepModel)
potrzebuje aż \resBestDeepTrainTime~s --- czyli kilkadziesiąt razy dłużej --- by
osiągnąć praktycznie ten sam wynik. Najbardziej rozbudowana sieć LSTM trenuje
się ponad osiem minut, a mimo to wypada słabiej. Również w fazie działania
(klasyfikacji nowego ruchu) las losowy i proste modele są wyraźnie szybsze od
sieci CNN i LSTM. Ten kompromis między kosztem a skutecznością dobrze obrazuje
rys.~\ref{fig:tradeoff}: najbardziej opłacalne modele leżą w lewej górnej części
wykresu --- wysoka skuteczność przy krótkim czasie uczenia.

\begin{figure}[H]
\centering
\includegraphics[width=0.8\textwidth]{time_vs_f1.png}
\caption{Kompromis między czasem treningu (skala logarytmiczna) a skutecznością
($F_1$). Modele najbardziej efektywne znajdują się w lewym górnym obszarze.}
\label{fig:tradeoff}
\end{figure}

\subsection{Przebieg uczenia modeli głębokich}

Rys.~\ref{fig:learning} przedstawia krzywe uczenia modeli głębokich: stratę
walidacyjną oraz ROC-AUC na zbiorze walidacyjnym w kolejnych epokach. Widoczna
jest stabilna zbieżność i działanie mechanizmu wczesnego zatrzymania, który
zapobiega przeuczeniu.

\begin{figure}[H]
\centering
\includegraphics[width=0.88\textwidth]{learning_curves.png}
\caption{Krzywe uczenia modeli głębokich (strata i ROC-AUC walidacyjny).}
\label{fig:learning}
\end{figure}

\subsection{Interpretowalność modeli i ważność cech}

Dodatkową zaletą lasu losowego jest interpretowalność. Rys.~\ref{fig:importance}
przedstawia najważniejsze cechy według miary ważności Giniego. Pozwala to
zidentyfikować, które statystyki przepływu najsilniej różnicują ruch złośliwy i
normalny --- informacja cenna dla analityków bezpieczeństwa, a trudniej dostępna
w modelach głębokich, działających jak ,,czarna skrzynka''.

\begin{figure}[H]
\centering
\includegraphics[width=0.8\textwidth]{feature_importance.png}
\caption{Najważniejsze cechy według lasu losowego (ważność Giniego).}
\label{fig:importance}
\end{figure}

\subsection{Podsumowanie wyników}

Zebrane wyniki pozwalają wstępnie odpowiedzieć na postawione pytania:
\begin{itemize}
  \item \textbf{Skuteczność:} las losowy i najlepszy model głęboki
        (\resBestDeepModel) są praktycznie nierozróżnialne pod względem $F_1$ i
        ROC-AUC; las losowy dodatkowo wykrywa najwięcej ataków (najwyższa
        czułość).
  \item \textbf{Koszt:} las losowy jest zdecydowanie tańszy --- uczy się
        kilkadziesiąt razy szybciej niż sieci głębokie, oferując najlepszą
        relację skuteczność--koszt.
  \item \textbf{Złożoność:} rozbudowane sieci sekwencyjne (zwłaszcza LSTM) nie
        przyniosły przewagi --- większa złożoność nie przełożyła się na lepsze
        wyniki, co potwierdza obserwacje z literatury dotyczące danych
        tabelarycznych.
\end{itemize}
